% ---------- PROFESSIONAL EXPERIENCE ----------
% Confirm client/data volumes and any product-specific performance figures before adding metrics.
\section*{Professional Experience}

\resumeEntry{Data Scientist \& Machine Learning Engineer}{Feb 2026 -- Present}{Deemea}{Paris, France}
\begin{resumeItemize}
  \resumeItem{Designed automated \textbf{DICOM data ingestion and processing pipelines} integrating hospital PACS and cloud platforms to support multi-center clinical studies.}
  \resumeItem{Developed deep learning-based \textbf{image quality control (QC)} algorithms to detect motion artifacts, slice mismatches, and corrupted data before clinical review.}
  \resumeItem{Developed computer vision algorithms for \textbf{diagnostic support, segmentation, and classification} of CT, X-ray, and MRI scans across orthopedics, oncology, and rheumatology.}
  \resumeItem{Built and deployed containerized AI services using \textbf{Docker and PyTorch}, integrating automated pre-annotations and interactive 3D visualizations into clinical workflows.}
  \resumeItem{Implemented model evaluation and \textbf{MLOps workflows} covering calibration and data-drift analysis to assess robustness across multi-site hospital datasets.}
  \resumeItem{Worked with radiologists, software engineers, and clinical stakeholders to translate clinical requirements into production-ready AI features.}
\end{resumeItemize}

\vspace{3pt}
\resumeEntry{Ph.D. Researcher in Machine Learning \& Computational Pathology}{Jan 2022 -- Jan 2026}{CentraleSupélec (OPIS / CVN Lab) \& AP-HP Hôpital Bicêtre}{Gif-sur-Yvette, France}
\begin{resumeItemize}
  \resumeItem{\textbf{Thesis:} \href{https://aymensadraoui.github.io/}{\textit{AI-Driven Methods in Histopathology for the Diagnosis and Prognosis of Hepatocellular Carcinoma}}. Developed deep learning models for histopathological architecture classification, hepatocellular carcinoma diagnosis, and early recurrence prediction using whole-slide images (WSIs) and clinical metadata.}
  \resumeItem{Designed \textbf{unrolled proximal optimization networks} (AlterPGA, U-AlterPGA, and U-AlterFISTA) for stain separation and normalization of digitized histopathology slides.}
  \resumeItem{Developed optimization-based methods for \textbf{transductive few-shot learning}, including UNEM and PADDLE-based approaches; U-PADDLE achieved 80.3\% accuracy in the evaluated 100-shot benchmark.}
  \resumeItem{Built scalable \textbf{ETL and image-processing pipelines} using PyTorch, OpenCV, Scikit-learn, OpenSlide, and Optuna to ingest, tile, and curate more than 4 TB of gigapixel whole-slide image data.}
  \resumeItem{Designed multi-center evaluation frameworks addressing \textbf{generalization, calibration, uncertainty, and interpretability}; the HCC architecture classifier achieved 82.3\% accuracy on the external Henri-Mondor cohort.}
  \resumeItem{Collaborated with pathologists and clinical researchers to translate histopathological criteria into quantitative machine learning objectives and clinically relevant predictive features.}
\end{resumeItemize}
